\newsection{Mathematische Konzepte}{task-3}

\begin{enumerate}
  \item {
    Sei $S = \{2, 5, 8, 11, 14\}$ eine Menge.
    Zeigen Sie mittels Diagonalisierung, dass die Menge \\
    $M = \{f~|~f: \mathbb{N} \to S\}$ nicht abzählbar unendlich ist, ohne die
    Überabzählbarkeit anderer Mengen zu nutzen.
  }
  \newpage
  \item {
    Bestimmen Sie mit Hilfe des chinesischen Restsatzes eine Zahl
    $x \in \mathbb{Z}$ mit $x \equiv 3$ mod $12$ und $x \equiv 4$ mod $7$.
    Geben Sie dazu ihre Zwischenschritte und das finale Ergebnis an.
  }
  \vfill
  \item {
    Wie viele verschiedene Worte kann man durch Permutation des Wortes WASSEREIS
    bilden?
    Begründen Sie ihre Rechnung kurz!
    Sie brauchen das Ergbnis nicht auszurechnen.
  }
  \vfill
  \newpage
  \item {
    Sei $G \coloneq \text{Abb}(\mathbb{R})$ die Menge aller Funktionen
    $f: \mathbb{N} \to \mathbb{N}$.
    Im Folgenden seien $f_1$ und $f_2$ zwei beliebige Funktionen aus $G$ und
    $n \in \mathbb{N}$ beliebig.
    Wir betrachten die folgende Verknüpfung auf diesre Menge:

    \[
      (G, \sim) \text{~mit~} (f_1 \sim f_2)(n) = f_1(n) \cdot f_2(n).
    \]

    Geben Sie für diese Struktur an, ob es sich um eine Halbgruppe, ein Monoid,
    eine Gruppe oder keines von alledem handelt und ob die Verknüpfung kommutativ
    ist.
    Begründen Sie Ihre Antworten.
  }
\end{enumerate}